\chapter{ROS 2 のインストールと環境構築}
\label{chap:ros2_install}

% この章で学ぶこと
\begin{goalbox}
  \begin{itemize}
    \item Ubuntu 24.04 に ROS 2 Jazzy をインストールする方法
    \item \cmd{source /opt/ros/jazzy/setup.bash} の意味と、\cmd{.bashrc} への設定
    \item turtlesim を使った ROS 2 の動作確認
    \item \cmd{ROS\_DOMAIN\_ID} を使った通信の分け方と、ネットワークの設定
  \end{itemize}
\end{goalbox}

この章では、Ubuntu 24.04 に ROS 2 Jazzy をインストールし、実際に動かしてみます。
ここまでの章で学んできた、パッケージ管理（\ref{chap:package}章）、環境変数と \cmd{source}（\ref{chap:env}章）、ネットワーク（\ref{chap:network}章）の知識が、さっそく役に立ちます。

\section{ROS 2 Jazzy のインストール}
\label{sec:ros2-install-install}

ROS 2 は、Ubuntu の公式リポジトリではなく、ROS 2 の開発元が用意しているリポジトリから、apt でインストールします（\ref{sec:package-repository}節）。
手順は変わることがあるので、実際にインストールするときは、公式ドキュメント（\url{https://docs.ros.org/en/jazzy/}）の「Installation」の「Ubuntu (deb packages)」のページも確認してください。
本書の執筆時点では、次のような手順です。

\subsection{文字コードの設定を確認する}

ROS 2 は、文字コードが UTF-8 の環境で動かすことを前提にしています。
\cmd{locale} コマンドで、今の設定を確かめましょう。

\begin{terminal}[文字コードの設定を確認する]
$ locale
LANG=ja_JP.UTF-8
LANGUAGE=ja:en
LC_CTYPE="ja_JP.UTF-8"
...
\end{terminal}

\cmd{UTF-8} が含まれていれば、そのままで構いません。
Ubuntu を日本語や英語で普通にインストールした場合は、UTF-8 になっているはずです。

\subsection{ROS 2 のリポジトリを追加する}

まず、Ubuntu の \cmd{universe} リポジトリ（\ref{sec:package-repository}節）が有効になっていることを確かめます。

\begin{terminal}[universe リポジトリを有効にする]
$ sudo apt install software-properties-common
$ sudo add-apt-repository universe
\end{terminal}

次に、ROS 2 のリポジトリの設定と GPG 鍵をまとめてインストールしてくれる \cmd{ros2-apt-source} というパッケージをダウンロードして、インストールします。

\begin{terminal}[ROS 2 のリポジトリを追加する]
$ sudo apt update && sudo apt install curl -y
$ export ROS_APT_SOURCE_VERSION=$(curl -s https://api.github.com/repos/ros-infrastructure/ros-apt-source/releases/latest | grep -F "tag_name" | awk -F\" '{print $4}')
$ curl -L -o /tmp/ros2-apt-source.deb "https://github.com/ros-infrastructure/ros-apt-source/releases/download/${ROS_APT_SOURCE_VERSION}/ros2-apt-source_${ROS_APT_SOURCE_VERSION}.$(. /etc/os-release && echo $VERSION_CODENAME)_all.deb"
$ sudo dpkg -i /tmp/ros2-apt-source.deb
\end{terminal}

長いコマンドですが、これまでの章で学んだことの組み合わせです。

\begin{itemize}
  \item 2 行目では、\cmd{curl}（Web からデータを取得するコマンド）で \cmd{ros2-apt-source} の最新版の情報を取得し、パイプ（\ref{sec:text-pipe-pipe}節）でつないだ \cmd{grep} と \cmd{awk}（\ref{sec:text-pipe-tools}節）でバージョン番号だけを取り出して、環境変数 \cmd{ROS\_APT\_SOURCE\_VERSION} に入れています。\cmd{\$(コマンド)} は、コマンドの実行結果を使う書き方です（\ref{sec:shellscript-variable}節）。
  \item 3 行目では、そのバージョンの \cmd{.deb} ファイルを \cmd{/tmp} にダウンロードしています。\cmd{\$(. /etc/os-release \&\& echo \$VERSION\_CODENAME)} の部分は、Ubuntu のコードネーム（\cmd{noble}）に置き換わります。
  \item 4 行目の \cmd{dpkg -i} は、ダウンロードした \cmd{.deb} ファイルをインストールするコマンドです。これで、ROS 2 のリポジトリの設定ファイルと GPG 鍵が登録されます。
\end{itemize}

\subsection{ROS 2 をインストールする}

リポジトリを追加したら、パッケージの一覧を更新し、インストール済みのパッケージを最新にしてから、ROS 2 をインストールします。

\begin{terminal}[ROS 2 Jazzy をインストールする]
$ sudo apt update
$ sudo apt upgrade
$ sudo apt install ros-jazzy-desktop
$ sudo apt install ros-dev-tools
\end{terminal}

\cmd{ros-jazzy-desktop} は、ROS 2 の本体に加えて、RViz などの GUI のツールや、デモ・チュートリアル用のパッケージをまとめたものです。
ダウンロードするデータが大きいので、インストールには時間がかかります。
\cmd{ros-dev-tools} には、パッケージをビルドするための \cmd{colcon} や、依存関係を解決する \cmd{rosdep}（\ref{chap:workspace}章）など、開発に使う道具が含まれています。

\begin{tablebox}[ROS 2 Jazzy の主なインストール用パッケージ][tab:ros2-install-packages]
  \begin{tabular}{lp{0.55\linewidth}}
    \toprule
    パッケージ & 内容 \\
    \midrule
    \cmd{ros-jazzy-desktop} & ROS 2 本体・GUI のツール（RViz など）・デモ。本書ではこれを使う \\
    \cmd{ros-jazzy-ros-base} & GUI のツールを含まない最小限の構成。ロボットに載せる PC など向け \\
    \cmd{ros-dev-tools} & colcon・rosdep などの開発用の道具 \\
    \bottomrule
  \end{tabular}
\end{tablebox}

ROS 2 の個々のパッケージは、\cmd{ros-jazzy-パッケージ名}（パッケージ名のアンダースコアはハイフンに変える）という名前で apt からインストールできます。
たとえば、\ref{chap:practice}章で使う \cmd{teleop\_twist\_keyboard} は \cmd{ros-jazzy-teleop-twist-keyboard} です。

\begin{point}[仮想マシンのスナップショット]
  ROS 2 のインストールが終わったら、VirtualBox のスナップショット（\ref{sec:linux-setup}節）をとっておきましょう。
  この後の章でいろいろ試して環境がおかしくなっても、この時点にすぐ戻れます。
\end{point}

\section{環境のセットアップ}
\label{sec:ros2-install-setup}

\subsection{ROS 2 の設定を読み込む}

ROS 2 は \cmd{/opt/ros/jazzy} にインストールされます（\ref{sec:linux-internals-fs}節）。
ところが、インストールしただけでは、まだ \cmd{ros2} コマンドは使えません。

\begin{terminal}[インストールしただけでは ros2 コマンドは使えない]
$ ros2 --help
ros2: command not found
\end{terminal}

\cmd{ros2} コマンドが置かれているディレクトリが、\cmd{PATH}（\ref{sec:env-variable}節）に含まれていないからです。
ROS 2 を使うには、ターミナルを開くたびに、次のコマンドで ROS 2 の設定を読み込む必要があります。

\begin{terminal}[ROS 2 の設定を読み込む]
$ source /opt/ros/jazzy/setup.bash
$ ros2 --help
usage: ros2 [-h] [--use-python-default-buffering] Call `ros2 <command> -h` for more detailed usage. ...
...
\end{terminal}

\cmd{setup.bash} は、ROS 2 を使うための環境変数を設定するシェルスクリプトです。
\ref{sec:env-export-source}節で学んだように、環境変数の設定を今のシェルに反映させるために、\cmd{bash} ではなく \cmd{source} で読み込みます。
読み込んだ後は、ROS 2 に関係する環境変数が設定されていることを確かめられます。

\begin{terminal}[ROS 2 の環境変数を確認する]
$ printenv | grep -i ros
ROS_VERSION=2
ROS_PYTHON_VERSION=3
ROS_DISTRO=jazzy
...
\end{terminal}

\subsection{\cmd{.bashrc} に設定を追加する}

毎回 \cmd{source /opt/ros/jazzy/setup.bash} を入力するのは面倒なので、\cmd{.bashrc}（\ref{sec:env-bashrc}節）に追加して、ターミナルを開くたびに自動で読み込まれるようにしておきます。
同じ行が何度も追加されないように、\cmd{grep} ですでに書かれていないかを確かめてから追加します。

\begin{terminal}[.bashrc に ROS 2 の設定を追加する]
$ grep "source /opt/ros/jazzy/setup.bash" ~/.bashrc
$ echo "source /opt/ros/jazzy/setup.bash" >> ~/.bashrc
$ source ~/.bashrc
\end{terminal}

1 行目で何も表示されなければ、まだ書かれていないので、2 行目で追加します（\cmd{>>} を \cmd{>} と間違えないように注意しましょう）。
テキストエディター（\ref{sec:files-edit}節）で \cmd{\textasciitilde/.bashrc} を開いて、末尾に書き足しても構いません。
新しく開いたターミナルで \cmd{ros2 --help} が実行できれば、設定は完了です。

\begin{point}[ros2 コマンドの補完]
  ROS 2 の設定を読み込むと、\cmd{ros2} コマンドでも Tab 補完（\ref{sec:terminal-keys}節）が使えるようになります。
  \cmd{ros2 r} まで入力して \key{Tab} を押すと \cmd{ros2 run} に、さらにパッケージ名やノード名の途中まで入力して \key{Tab} を押すと、残りが補完されます。
  ROS 2 のコマンドは長いものが多いので、積極的に使いましょう。
\end{point}

\section{turtlesim で動作確認}
\label{sec:ros2-install-turtlesim}

\subsection{turtlesim を起動する}

ROS 2 が正しくインストールできたかを、\textbf{turtlesim} というシミュレータで確かめましょう。
turtlesim は、画面の中のカメを ROS 2 で動かす、学習用のシンプルなシミュレータです。
\cmd{ros-jazzy-desktop} に含まれていますが、入っていない場合は \cmd{sudo apt install ros-jazzy-turtlesim} でインストールします。

ノードを起動するには、\cmd{ros2 run パッケージ名 実行ファイル名} を使います。
次のコマンドで、turtlesim のノードを起動します。

\begin{terminal}[turtlesim のノードを起動する]
$ ros2 run turtlesim turtlesim_node
[INFO] [1727930000.123456789] [turtlesim]: Starting turtlesim with node name /turtlesim
[INFO] [1727930000.234567890] [turtlesim]: Spawning turtle [turtle1] at x=[5.544445], y=[5.544445], theta=[0.000000]
\end{terminal}

青い背景の中にカメが 1 匹いるウィンドウが開けば成功です（図\ref{fig:ros2-install-turtlesim}）。
ノードは止めるまで動き続けるので、このターミナルはそのままにしておきます。

\begin{figure}[tb]
  \centering
  \large
  \begin{tikzpicture}
    \draw rectangle++(4.5,3);
    \clip rectangle++(4.5,3);
    \foreach\x in{0,1,2,3,4}{
      \foreach\y in{0,1,2,3,4,5,6}{
        \path(\x*1.5-\y*0.5,\y*0.6)node[rotate=10,text=gray]{\sffamily\bfseries\itshape Sample };
      }
    }
  \end{tikzpicture}
  %% 入れる画像：turtlesim を起動し、turtle_teleop_key でカメを少し動かして軌跡が描かれた状態のウィンドウのスクリーンショット
  % eps 画像を貼る場合は includegraphics をお使いください。
  % \includegraphics[width=0.6\linewidth]{ros2_install/turtlesim.png}
  \caption{turtlesim のウィンドウ}
  \label{fig:ros2-install-turtlesim}
\end{figure}

\subsection{キーボードでカメを動かす}

\key{Ctrl}+\key{Shift}+\key{T} で、ターミナルに新しいタブを開きます（\key{Ctrl}+\key{Alt}+\key{T} で新しいウィンドウを開いても構いません）。
新しいタブで、キーボードの入力を受け付けるノードを起動します。

\begin{terminal}[キーボードで操作するノードを起動する]
$ ros2 run turtlesim turtle_teleop_key
Reading from keyboard
---------------------------
Use arrow keys to move the turtle.
...
'Q' to quit.
\end{terminal}

このターミナルを選んだ状態で矢印キーを押すと、turtlesim のウィンドウの中のカメが動き、通った跡に線が描かれます。
カメが動かないときは、このターミナルのウィンドウがアクティブになっている（クリックして選ばれている）かを確かめてください。

\subsection{何が起きているのか}

このとき、2 つのノードが、ROS 2 の通信を使ってデータをやり取りしています（図\ref{fig:ros2-install-graph}）。

\begin{figure}[tb]
  \centering
  % 図：turtlesim のノードとトピック（chapters/ros2/ros2_install.tex）
\begin{tikzpicture}[
    node/.style={draw, ellipse, align=center, font=\small, fill=blue!8, minimum height=10mm, inner sep=2pt},
    topic/.style={draw, rectangle, font=\footnotesize\ttfamily, fill=orange!10, inner sep=3pt}]
  \node[node] (teleop) at (0,0) {\texttt{/teleop\_turtle}\\{\footnotesize キー入力}};
  \node[topic] (cmd) at (4.9,0) {/turtle1/cmd\_vel};
  \node[node] (sim) at (9.2,0) {\texttt{/turtlesim}\\{\footnotesize カメの画面}};
  \draw[figarrow] (teleop) -- node[figlabel, above]{速度の指令} (cmd);
  \draw[figarrow] (cmd) -- (sim);
  \node[figlabel, text=black!60, below=3mm of teleop] {ノード（楕円）};
  \node[figlabel, text=black!60, below=3mm of cmd] {トピック（四角）};
\end{tikzpicture}

  \caption{turtlesim を動かしているときのノードとトピック}
  \label{fig:ros2-install-graph}
\end{figure}

キーボードの入力を受け付ける \cmd{/teleop\_turtle} というノードが、押されたキーに応じた速度の指令を、\cmd{/turtle1/cmd\_vel} という名前の\textbf{トピック}に送っています。
カメの画面を表示する \cmd{/turtlesim} というノードが、そのトピックから指令を受け取って、カメを動かしています。
さらに別のタブを開いて、動いているノードとトピックの一覧を確かめてみましょう。

\begin{terminal}[ノードとトピックの一覧を確認する]
$ ros2 node list
/teleop_turtle
/turtlesim
$ ros2 topic list
/parameter_events
/rosout
/turtle1/cmd_vel
/turtle1/color_sensor
/turtle1/pose
\end{terminal}

ノードやトピックについては、次の\ref{chap:ros2_concepts}章で詳しく学びます。
確認が終わったら、それぞれのターミナルで \key{Ctrl}+\key{C} を押してノードを止めます（\cmd{turtle\_teleop\_key} は \key{q} でも終了できます）。

\subsection{デモのノードで通信を確かめる}

もう 1 つ、メッセージを送るノードと受け取るノードのデモでも試してみましょう。
2 つのターミナルで、それぞれ次のコマンドを実行します。

\begin{terminal}[1 つ目のターミナル：メッセージを送るノード]
$ ros2 run demo_nodes_cpp talker
[INFO] [1727930100.123456789] [talker]: Publishing: 'Hello World: 1'
[INFO] [1727930101.123456789] [talker]: Publishing: 'Hello World: 2'
...
\end{terminal}

\begin{terminal}[2 つ目のターミナル：メッセージを受け取るノード]
$ ros2 run demo_nodes_py listener
[INFO] [1727930102.234567890] [listener]: I heard: [Hello World: 3]
...
\end{terminal}

\cmd{talker} は C++ で、\cmd{listener} は Python で書かれたノードです。
書かれた言語が違っても、ROS 2 のノードどうしは同じように通信できることがわかります（\ref{sec:python-pros-cons}節のポイント）。

\section{ROS\_DOMAIN\_ID とネットワーク設定}
\label{sec:ros2-install-network}

\subsection{同じネットワークのノードは自動でつながる}

\ref{sec:what-is-ros-architecture}節で説明したように、ROS 2 では、同じネットワークにつながっているコンピュータのノードどうしが、自動でお互いを見つけて通信します。
これは便利な反面、研究室や競技会場のように、同じネットワークに複数の人のロボットや PC がある環境では困ったことになります。
自分のノードが、隣の人のロボットのトピックを受け取ってしまったり、自分の指令で隣の人のロボットが動いてしまったりするからです。

\subsection{ROS\_DOMAIN\_ID で通信を分ける}

このような混乱を防ぐのが、\textbf{\cmd{ROS\_DOMAIN\_ID}} という環境変数です。
\cmd{ROS\_DOMAIN\_ID} には 0 から 101 までの番号を設定でき、\textbf{同じ番号のノードどうしだけが通信します}。
番号を設定しなかった場合は 0 が使われます。

\begin{terminal}[ROS\_DOMAIN\_ID を設定する]
$ export ROS_DOMAIN_ID=12
\end{terminal}

\cmd{export} で設定した値は、そのターミナルだけで有効です（\ref{sec:env-export-source}節）。
いつも同じ番号を使うなら、\cmd{.bashrc} に追加しておきます。

\begin{terminal}[.bashrc に ROS\_DOMAIN\_ID を追加する]
$ echo "export ROS_DOMAIN_ID=12" >> ~/.bashrc
\end{terminal}

\begin{point}[ROS\_DOMAIN\_ID の決め方]
  研究室やチームの中で、人やロボットごとに重ならない番号を割り当てておきましょう。
  開発用 PC とロボットの PC で通信させたいときは、両方に同じ番号を設定します。
  「ノードが見えない」「トピックが届かない」というトラブルの原因として、\cmd{ROS\_DOMAIN\_ID} の設定の食い違いはとてもよくあります（\ref{chap:debug}章）。
  うまく通信できないときは、まず両方のターミナルで \cmd{echo \$ROS\_DOMAIN\_ID} を実行して、番号が同じかを確かめましょう。
\end{point}

\subsection{自分のコンピュータの中だけで通信する}

ほかのコンピュータとは一切通信せず、自分のコンピュータの中のノードどうしだけで通信したいときは、\cmd{ROS\_AUTOMATIC\_DISCOVERY\_RANGE} という環境変数に \cmd{LOCALHOST} を設定します。
シミュレータだけで学習しているときなどに便利です。

\begin{terminal}[自分のコンピュータの中だけで通信する]
$ export ROS_AUTOMATIC_DISCOVERY_RANGE=LOCALHOST
\end{terminal}

古い資料には \cmd{ROS\_LOCALHOST\_ONLY=1} という設定が書かれていることがありますが、ROS 2 Jazzy では \cmd{ROS\_AUTOMATIC\_DISCOVERY\_RANGE} を使います。

\subsection{別のコンピュータと通信するときの注意}

開発用 PC とロボットの PC のように、別々のコンピュータのノードどうしを通信させるには、次の点を確かめておきましょう。

\begin{itemize}
  \item 2 台のコンピュータが同じネットワークにつながっていて、\cmd{ping}（\ref{sec:network-check}節）で通信できること
  \item 両方の \cmd{ROS\_DOMAIN\_ID} が同じであること
  \item VirtualBox の仮想マシンを使っている場合は、ネットワークの設定が「ブリッジアダプター」になっていること（\ref{sec:linux-setup}節）。標準の NAT のままでは、ほかのコンピュータのノードを見つけられません
\end{itemize}

準備ができたら、片方のコンピュータで \cmd{talker}、もう片方で \cmd{listener} を起動して、メッセージが届くかを確かめてみましょう。
届かない場合は、ネットワークがノードを見つけるための通信（マルチキャスト）に対応しているかを、\cmd{ros2 multicast} コマンドで確かめられます。

\begin{terminal}[2 台のコンピュータの間でマルチキャストを確かめる]
$ ros2 multicast receive          # 1 台目：受信を待つ
$ ros2 multicast send             # 2 台目：送信する
\end{terminal}

1 台目に \cmd{Received from ...} と表示されれば、マルチキャストで通信できています。
Wi-Fi のルータによっては、マルチキャストの通信がうまく届かないことがあります。

\section*{この章のまとめ}

\begin{itemize}
  \item ROS 2 Jazzy は、ROS 2 のリポジトリを追加してから、apt で \cmd{ros-jazzy-desktop} と \cmd{ros-dev-tools} をインストールします。手順は公式ドキュメントも確認しましょう。
  \item ROS 2 を使うには、\cmd{source /opt/ros/jazzy/setup.bash} で設定を読み込みます。\cmd{.bashrc} に書いておくと、ターミナルを開くたびに自動で読み込まれます。
  \item \cmd{ros2 run パッケージ名 実行ファイル名} でノードを起動します。turtlesim で、ノードどうしがトピックで通信する様子を確かめました。
  \item 同じ \cmd{ROS\_DOMAIN\_ID} のノードどうしだけが通信します。人やロボットごとに番号を分け、通信させたいコンピュータどうしは同じ番号にそろえます。
\end{itemize}
